% New independent proof audit. No documentclass; no old file is edited.
\subsection{From the actual finite upper tree to bounded convolution atoms}
\label{sec:actual-k-compilation-joint-audit}

This audit proves the finite combinatorial and factorization interface
used by \texttt{lower-k-uniform-product-interface-review.tex}. The
objects being compiled are actual signed sieve coefficients. Merely
factoring each integer in their support would not suffice. The
dispersion theorem, outer-product HB/SW argument and limiting density
evaluation remain separate inputs. The finite tests recorded in the
companion JSON are diagnostics of the identities below, not their proof.

\paragraph{1. The raw tree and the sign of its first exits.}
At any fixed $N$ order the allowed sifting primes decreasingly. A
history contains the ordered prefix, its coarse labels, inherited
admission decisions, selected mask and frozen level. Let $\mathcal T$
contain the empty tuple, be closed under deleting final slots, and have
only odd-length first excluded prefixes. Put
$\lambda(d)=\mu(d)\mathbf1_{d\in\mathcal T}$ for squarefree $d$.
For every finite set $Q$ of allowed primes and $0\le h(p)\le1$,
\begin{equation}\label{eq:actual-k-first-exit}
 \sum_{d\mid Q}\lambda(d)h(d)
 =\prod_{p\mid Q}(1-h(p))+
 \sum_{P=p_1\cdots p_j\mid Q\atop P\text{ first exit}}
 h(P)\prod_{p\mid Q,\ p<p_j}(1-h(p)).
\end{equation}
Indeed, partition every excluded divisor by its unique first excluded
prefix $P$. Its remaining factors are an arbitrary subset of the
primes below $p_j$. Its contribution to the full Euler product is
$(-1)^jh(P)\prod_{p<p_j}(1-h(p))$, and $j$ is odd. Subtracting these
contributions proves the identity. This proof permits history-dependent
levels; the history of a prefix must be identical in all descendants.

In the inspected tree the first stop is at rank one; the old third
decision is at rank three. A fourth upgrade preserves its rank-four
node and changes subsequent odd tests. A fifth upgrade replaces the
fifth test itself, leaving the earlier decisions intact. A sixth
upgrade preserves membership of its rank-six node and changes only
tests $7,9,\ldots$. Reverting an upgrade on a boundary strip means
restoring the inherited rule there. These operations preserve the
hypotheses of \eqref{eq:actual-k-first-exit}. They do not mean that
signed finite weights are monotone in the chosen level.

In particular an extra exclusion at rank six is forbidden. The need is
substantive: on two ordered primes, retaining $1,p_1,p_2$ but excluding
$p_1p_2$ gives divisor sum $1-1-1=-1$ at $n=p_1p_2$, below the required
zero. The actual rank-six construction avoids this operation.

\paragraph{2. Rough projection has the required one-sided density.}
Let $P_N(y)$ use only primes not dividing $N$, let
$g_N(p)=1/(p-1)$ on these primes, and let
$V_N(y)=\prod_{p<y,p\nmid N}(1-g_N(p))$. Since $N$ is even,
$0\le g_N(p)<1$. Fix $w<z$ and restrict the same raw history rule
to $Q=P_N(z)/P_N(w)$. Multiplication of
\eqref{eq:actual-k-first-exit} by $V_N(w)$ gives
\[
 V_N(w)\sum_{d\mid Q}\lambda(d)g_N(d)
 =V_N(z)+\sum_{P\text{ full first exit}\atop p_j\ge w}
                g_N(P)V_N(p_j)
 \le \sum_{d\mid P_N(z)}\lambda(d)g_N(d).
\]
A rough first exit is exactly a full first exit whose last prime is
at least $w$, because all its earlier primes are larger. The missing
terms are nonnegative. Thus the rough density times $V_N(w)$ is no
larger than the full density. It is not asserted equal to it.
If the raw definition deactivates deep upgrades below a cutoff, use
that same definition on both sides; passage to an endpoint tree is a
separate limiting-density argument.

\paragraph{3. The correct box inequality, including endpoints.}
Choose a finite fine partition into disjoint half-open prime intervals,
refining every coarse label. Keep each coarse history decision fixed
on its labelled box. At an even final rank use lower endpoints in all
applicable historical odd cube tests and non-strict inequalities. At
an odd final rank use upper endpoints and strict inequalities. Define
$c(d)=(-1)^m$ on retained squarefree rough products of rank $m$, and
zero otherwise. Take the rank cutoff larger than every possible raw
rough rank; this is possible since $d\le N^{\theta_*}$ and
$p\ge N^\omega$ with fixed $\omega>0$.

A positive raw coefficient always passes the even outer tests. Every
negative boxed coefficient passes the actual strict odd tests. Hence
\[
 c(d)\ge\lambda(d)\quad\text{coefficientwise},\qquad
 \sum_{d\mid n,\ d\mid Q}c(d)\ge\mathbf1_{(n,Q)=1}.
\]
The boxed support itself need not be prefix closed. Equality at a
grid endpoint belongs to exactly one half-open interval; the stated
strict/non-strict tests give these pointwise conclusions even there.
Every applicable historical cube must be retained. Testing only a
new final cap would incorrectly forget an earlier stop.

If the box width in exponent coordinates is $\tau$, a discrepancy
lies within $O(M\tau)$ of a historical odd cube at one of the finitely
many frozen levels. Fix all primes except the tagged prime in this
cube. Its reciprocal-prime mass is $O(M\tau/\omega)+o_N(1)$.
The remaining subsets sum to at most
$\prod_{w\le p<z}(1+g_N(p))\ll\omega^{-1}$.
There are $O(M)$ possible tags, $M=O(\omega^{-1})$, and
$V_N(w)/V_N(z)\ll\omega^{-1}$. Thus, with a constant depending on
the frozen cap list,
\[
 0\le V_N(w)\sum_d(c(d)-\lambda(d))g_N(d)
 \ll V_N(z)C_{\mathcal T}\tau\omega^{-5}+o_N(V_N(z)).
\]
This estimate also covers a tag among the marked prefix primes,
because one sums over the finite frozen levels rather than
differentiating a discontinuous cap. Gate strips are first treated by
inherited legal rules, not by arbitrary even-rank deletion.

\paragraph{4. Sign-sensitive small-prime attachment.}
Let $\beta^\pm$ be squarefree fundamental-lemma weights supported on
primes below $w$, with $|\beta^\pm(b)|\le1$, and put
\[
 B^-(n)=\sum_{b\mid n}\beta^-(b)
 \le I_w(n):=\mathbf1_{(n,P_N(w))=1}
 \le B^+(n)=\sum_{b\mid n}\beta^+(b).
\]
Attach $\beta^+$ to $c(d)>0$, and $\beta^-$ to $c(d)<0$.
For every individual rough divisor,
\[
 c(d)B^{\operatorname{sgn}c(d)}(n)\ge c(d)I_w(n).
\]
Summing and using the previous paragraph proves
\[
 \sum_{q\mid n}\rho(q)\ge I_w(n)\sum_{d\mid n,d\mid Q}c(d)
 \ge\mathbf1_{(n,P_N(z))=1}.
\]
No positivity of $B^-(n)$ is used. Since the small/rough decomposition
$q=bd$ is unique, $|\rho(q)|\le1$. Repeated appearances in the
subsequent atom expansion do not alter this original coefficient norm.

If both beta densities differ from $V_N(w)$ by at most
$\varepsilon_\beta V_N(w)$, their attachment costs at most
\[
 \varepsilon_\beta V_N(w)\sum_d|c(d)|g_N(d)
 \ll \varepsilon_\beta\omega^{-2}V_N(z).
\]
Thus the density error is controlled even though negative coefficients
receive a different beta sequence. Combined with the rough comparison
and box error, this gives precisely an upper bound by the raw density
plus these two errors.

\paragraph{5. The exact atom identity and factorial bookkeeping.}
Fix a retained ordered fine-label vector $J=(J_1,\ldots,J_m)$.
For a label $I$, let $m_I$ be its multiplicity and let
$M_J=\prod_I m_I!$. Put
$T_J=\chi_{J_1}*\cdots*\chi_{J_m}$, with each prime indicator also
restricted to $p\nmid N$. If $A_J$ denotes the indicator of the sorted
squarefree products in this label box, unique factorization gives
\[
 \mathbf1_{\rm squarefree}(d)T_J(d)=M_J A_J(d).
\]
Every distinct prime product has exactly $m_I!$ assignments among the
slots of label $I$. Its rough contribution is therefore exactly
$(-1)^m\mathbf1_{\rm squarefree}T_J/M_J$, followed by the assigned
$\beta^+$ or $\beta^-$. In particular the reciprocal $M_J$ occurs
once, not once per new group and not twice.

Suppose a fixed partition $\mathcal I_1,\mathcal I_2,\mathcal I_3$
of its slots has been supplied. For $m_i=|\mathcal I_i|$, define
\[
 \gamma_i=\frac1{m_i!}\mathop{*}_{j\in\mathcal I_i}\chi_{J_j},
 \qquad \gamma_i=\delta_1\text{ if }m_i=0.
\]
Unique factorization bounds each unrestricted convolution coefficient
by $m_i!$, even with repeated primes. Hence $|\gamma_i|\le1$, and
the exact unrestricted identity is
\[
 T_J=(m_1!m_2!m_3!)\gamma_1*\gamma_2*\gamma_3.
\]
The new scalar is the old scalar times $m_1!m_2!m_3!$; its original
factor $1/M_J$ is retained. Attach beta to the third group. Since its
primes and the rough primes are disjoint,
\[
 |(\gamma_3*\beta)(n)|
 =|\gamma_3(n_{\rm rough})\beta(n_{\rm small})|\le1.
\]
There are finitely many labels and bounded rough rank at fixed
parameters, so the atom count and total absolute scalar mass are
bounded independently of $N$. Their constants need not be small.

Removing the squarefree restriction to use these exact convolutions
only adds moduli divisible by $p^2$ for a prime $p\ge w$. For a fixed
divisor order $C$ and $Q_*=N^{\theta_*}$,
\[
 \sum_{d\le Q_*,\ \exists p\ge w:p^2\mid d}
 \tau_C(d)\left(\frac Nd+1+\frac N{\varphi(d)}\right)
 \ll_C (N+Q_*)w^{-1}(\log N)^{C'}.
\]
To see this, write $d=p^2e$, use
$\tau_C(p^2e)\le\tau_C(p^2)\tau_C(e)$ and
$\varphi(p^2e)\ge\varphi(p^2)\varphi(e)$, sum standard divisor
and reciprocal-divisor bounds over $e$, then use
$\sum_{p\ge w}p^{-2}\ll w^{-1}$. This is a power saving.
It applies before HB expansion to the bounded positive outer tuple
weight. Thus the audit does not silently assume a pointwise bound
for divisor coefficients inside an arbitrary progression.

\paragraph{6. Selected gates imply every required residual square.}
Write $t_1\ge\cdots\ge t_m>0$, $S_j=\sum_{i\le j}t_i$, and
select $I\subseteq\{1,\ldots,r\}$. Let $s_I=\sum_{i\in I}t_i$.
For an unselected slot $j$, the needed square exponent is
\[
 s_I+\sum_{i<j,\ i\notin I}t_i+2t_j.
\]
For $j\le r$ this is exactly the gate $G_I$, individually for every
unselected slot; the maximum cannot be replaced by only the first
such gate. For $j>r$ it equals $S_{j-1}+2t_j$. An admitted odd cube
gives its own square because $2t_j\le3t_j$. For an even descendant,
\[
 S_{j-1}+2t_j\le S_{j-2}+3t_{j-1}.
\]
The right side is the preceding admitted odd cube. The initial full
squares are supplied by $W_r$. Consequently the complete descendant
has both ordinary well support and every selected-remaining square
at its frozen cap. At rank six the node already existed; its inherited
fifth cube establishes admission, while the new gate checks all six
initial squares. At later upgrades apply the same argument to the
last active selected prefix, preserving all earlier admission decisions.

\paragraph{7. Two-capacity factorization and the capacity below one.}
Use local scale $x$, spectral parameter $0\le\sigma\le7/32$,
$1/2\le\vartheta<5/8$, $D=x^{\vartheta-2\delta}$ and
$x^\delta\le A\le x^{1/3-\delta/2}$. Put
\[
 v=2\vartheta-1,\qquad
 u=\min\left\{\frac{2-(2+2\sigma)\vartheta}{3-4\sigma},
                 \frac{2-2\vartheta-\alpha\sigma}{3}\right\}.
\]
Suppose a selected product satisfies $x^v\le b_0\le x^u$ and all
remaining-square conditions are at most $D$. If $C=D/(Ab_0)\ge1$,
process the remaining slots in their inherited order. Before a prime
$q$ is assigned, the square condition is
$q^2\le(A/a)(C/c)$. If neither group could accept $q$, the reverse
strict inequality would hold. Therefore a greedy choice assigns
every slot with $a\le A$, $c\le C$ and $b=b_0$. Directly,
\[
 A^2bc^2\le D^2/b\le x^{1-4\delta},\quad
 x^{\alpha\sigma}Aab^5c^2\le x^{\alpha\sigma}D^2b^3
                          \le x^{2-4\delta},
\]
\[
 (Ac/b)^{2\sigma}Aab^5c^2
 \le D^{2+2\sigma}b^{3-4\sigma}
 \le x^{2-(4+4\sigma)\delta}.
\]

If $C<1$, do not replace it by one while retaining $b_0$. Set
$B=D/A<b_0\le x^u$ and reassign all original slots into capacities
$A,B$. Both are at least one; in fact
$B\ge x^{1/6-3\delta/2}$. The original full square conditions give
the same two-capacity induction. It yields $d=a'b'$, $a'\le A$,
$b'\le B$, $c'=1$. Now
\[
 A^2b'\le AD\le x^{23/24-5\delta/2}\le x^{1-3\delta},
\]
\[
 x^{\alpha\sigma}Aa'b'^5\le x^{\alpha\sigma}D^2B^3
 \le x^{2-4\delta},\qquad
 (A/b')^{2\sigma}Aa'b'^5
 \le D^{2+2\sigma}B^{3-4\sigma}
 \le x^{2-(4+4\sigma)\delta}.
\]
The last $b'$ need not remain in the critical interval. That interval
was only a sufficient construction device; the required conclusions
are the three displayed factorization monomials.

For an inherited strict $W$ branch the buffered three-factor repair
is also valid: start with capacities $A,x^v,D/(Ax^v)$ and only
$p_1,p_2$ assigned. Its hypotheses are $p_i\le x^v$,
$p_j^2\le x^{1-\vartheta-2\delta}$, and $p_3\le x^{u-v}$.
If the greedy process never fails, substituting the three capacities
proves the monomials, using $u\ge v$. If $p_k$ is the first failure,
$b=b_0p_k$ lies in $(x^v,x^u]$ and
$Abc_0<(a_0p_k)(b_0p_k)c_0\le D$.
Thus $D/(Ab)>c_0\ge1$ and the two-capacity completion applies.
This verifies the initialization and failure step rather than relying
on the inaccurate unbuffered initialization in the printed argument.

\paragraph{8. One slot partition works throughout the box at each scale.}
At a fixed Type II length $U$ use $A_0=Ux^{-\zeta}$ with
$x^{2\zeta}\le U\le x^{1/3+\zeta/2}$. Then
$x^\zeta\le A_0\le x^{1/3-\zeta/2}$.
Apply the preceding construction with $\delta=\zeta$ to the ordered
upper endpoints of the fine prime intervals, after buffering the
support and selected gates. The greedy proofs use only ordered
positive sizes; primality or distinctness of those endpoints is
unnecessary. Keep their slot assignment fixed.

Every actual group product is at most its endpoint group product.
This remains true when an interval label occurs several times. In
the selected construction test $D/(A_0b_+)$ at the upper selected
product: if it is below one, the full-vector repair works on the
whole box, regardless of which actual selected products would have
larger capacity. The beta factor is attached only afterwards, enlarging
the third length by at most $N^\ell$. The first length stays below
$A_0$. The coefficients still represent the same $\rho$; the
partition may depend on $U$ and the atom but never on the actual
prime tuple inside this box.

Before small-factor and endpoint losses, replacing $A_0$ by $U$
costs exponents $2\zeta$, $\zeta$, and $(1+2\sigma)\zeta$ in the
three monomials. The available source savings are $3\zeta$,
$3\zeta$, and $2\zeta$. The least remaining saving is
$(1-2\sigma)\zeta\ge9\zeta/16$. Choose all added losses below this
fixed quantity and $\epsilon_{\rm disp}\ll\zeta$. After dyadic
subdivision the resulting bounded group coefficients meet
\[
 Q_1<Ux^{-\epsilon_{\rm disp}},\quad
 U^2Q_2Q_3^2<x^{1-8\epsilon_{\rm disp}},
\]
\[
 |a_{\rm res}|^\sigma UQ_1Q_2^5Q_3^2<x^{2-15\epsilon_{\rm disp}},
 \quad
 (UQ_3/Q_2)^{2\sigma}UQ_1Q_2^5Q_3^2
             <x^{2-15\epsilon_{\rm disp}}.
\]
These are the actual modulus-factor conditions in L23 Proposition 5.2.
Every exponent of a $Q_i$ is positive after combining the last
expression; in particular that of $Q_2$ is $5-2\sigma>0$. Smaller
dyadic lengths therefore preserve the bounds. Units are handled by
the point mass at one. Logarithmically many dyadic pieces do not
change the original coefficient or its finite scalar normalization.

\paragraph{9. A noncircular parameter choice.}
First fix a strict rounding budget $K-K_{\rm end}>0$ and one finite
buffered tree whose density is below $K$ by a fixed amount. The
justification of that density approximation is an inherited input;
it is not proved by this compilation. Freeze its coarse histories,
gate margins $e>0$, active prefix cutoffs and maximum level below
$3/5$. Fix the outer factor cutoff $a>0$. Then choose
$\zeta\ll\min(e,a)$ with $12\zeta<a/2$, choose
$\epsilon_{\rm disp}\ll\zeta$, and $\kappa\ll\epsilon_{\rm disp}$.
Local coordinates and $\alpha=1+O(\kappa)$ then preserve all margins.

Choose a single small beta precision $\eta$ next, for example with
$\ell=\nu_0\eta$ and $\omega=\nu_0\eta^2$ for a fixed positive
reference level $\nu_0$. Impose $\ell\ll\epsilon_{\rm disp}$,
make $\omega$ smaller than every active marked cutoff, and make
$e^{-c/\eta}\omega^{-2}$ smaller than the density budget.
This is possible because exponential decay beats the fixed power of
$1/\eta$. If the raw density itself still depends on $w$, choose
$\eta$ also in its already proved small-cutoff limiting range.
The resulting rank bound $M$ is now fixed. Only then choose the fine
width $\tau$ so that $M\tau\ll\epsilon_{\rm disp}$ and
$C_{\mathcal T}\tau\omega^{-5}$ fits the remaining density budget.
Impose also that the complete endpoint level plus $\ell+O(M\tau)$
is below $3/5$; ranks at most two stay below
$2(19/96)+\ell+O(M\tau)<1/2$.

Thus beta level and rough cutoff cannot be varied independently while
claiming the displayed beta error. The preceding coupled choice
supplies the intended reading of the earlier shorter parameter paragraph.
All choices, factorials and numbers of atoms are fixed before
$N\to\infty$. Smoothing and endpoint limits are taken afterwards.
There is no demand for uniformity as $a\downarrow0$, as the number of
labels grows, or as the depth tends to infinity.

\paragraph{Conclusion and remaining inputs.}
Under the explicitly stated legal-history and buffered-gate hypotheses,
the actual compiled coefficients are pointwise upper-sieve coefficients
with norm at most one, with density bounded above by the raw density
plus controllable errors, and with an exact finite signed atom expansion.
Each atom has a fixed slot partition at every required Type II scale
into three bounded convolution factors meeting the source monomials.
No counterexample was found to this scoped interface. The coupled
parameter choice and preservation of all historical cube tests above
make explicit two bookkeeping requirements that must not be omitted.

This conclusion is not a proof of distribution for arbitrary
coefficients supported on the same integers. It still needs the cited
dispersion theorem, the separately audited outer HB/SW and long-variable
argument, moving-residue range, positive radial boundary passage, and
the numerical and analytic limiting-density calculation of the chosen
fourth/fifth/sixth tree. This audit does not rerun those integrals or
independently certify the whole value $K=6.7619$.
